\section*{Chapter \ref{ch:b1:price-theory-at-work} --- Price Theory at Work}
% Numbers recomputed in code/bachelor-1/14-price-theory-at-work/exercises/solutions.py (tested).

\begin{solution}{exo:b1:price-theory-at-work:full}
$4 + 2 \times 15 = 34$ dalens.
\end{solution}

\begin{solution}{exo:b1:price-theory-at-work:read}
$100/3.93 = 25.4$ years of rent in 1980 and $100/2.49 = 40.2$ in 2020. The real
rate fell by $1.97 - (-1.29) = 3.26$ points.
\end{solution}

\begin{solution}{exo:b1:price-theory-at-work:quota-rent}
$40 \times (14 - 10) = 160$.
\end{solution}

\begin{solution}{exo:b1:price-theory-at-work:distance}
$0.05 \times 300 = 15$ dalens a tonne: the gap can be at most 15, and it is
exactly 15 if grain flows between the markets.
\end{solution}

\begin{solution}{exo:b1:price-theory-at-work:ceiling}
At 20 sellers offer 20 and buyers want 60: a shortage of 40. With the units
going to the buyers who value them most, the loss is the triangle between
$Q = 20$ and $Q = 40$, of height $60 - 20 = 40$ at $Q = 20$: $\tfrac12 \times 20 \times 40 = 400$. With
queues, buyers queue until the \omterm{def:b1:price-theory-at-work:full-price}{full price} is 60, the value of the last
unit served: $40 \times 20 = 800$ more is lost in time, 1,200 in all.
\end{solution}

\begin{solution}{exo:b1:price-theory-at-work:tariff}
Free trade: demand $120 - 50 = 70$, supply 20, imports 50. Tariff: demand $120 - 60 = 60$,
supply 24, imports 36; revenue $2 \times 36 = 72$. A quota of 36 needs a price at
which $120 - 5p - 2p = 36$, that is $p = 12$: the same. Consumer and \omterm{def:b1:partial-equilibrium-welfare:surplus}{producer
surplus} are the same in both; the rectangle of 72 is the state's revenue
under the tariff and the quota holders' rent under the quota.
\end{solution}

\begin{solution}{exo:b1:price-theory-at-work:path}
The price rises by $1.5\%$ a month: 100, 101.5, 103.02, 104.57, 106.14, 107.73,
109.34.
\end{solution}

\begin{solution}{exo:b1:price-theory-at-work:rate-fall}
$P/R = 1/(r + \delta)$: $1/0.06 = 16.67$ years of rent before and $1/0.04 = 25$ after. At an
unchanged rent the price rises by $25/16.67 - 1 = 50\%$.
\end{solution}

\begin{solution}{exo:b1:price-theory-at-work:corr}
The condition predicts a strong positive relation if depreciation and
expected growth are steady. Only the American numbers show it. In France
and Britain something else moved the yield: expected growth $g$, which
rises after years of rising prices and lowers the yield at a given real
rate; or the rent and price indices behind the yields.
\end{solution}

\begin{solution}{exo:b1:price-theory-at-work:derive}
Buy at $P$, rent out for $R$, pay upkeep $\delta P$, sell at $P(1 + g)$ at the end of the
year. The alternative is to lend $P$ at $r$. No profit either way requires
$R - \delta P + gP = rP$, that is $R = P(r + \delta - g)$. (Rent is taken as received at the
end of the year, and the rates as real.)
\end{solution}

\begin{solution}{exo:b1:price-theory-at-work:random}
Buyers' values run from 80 down to 20 over the 60 who want a unit; a random
buyer values a unit at $(80 + 20)/2 = 50$ on average, so the 20 units are worth
$20 \times 50 = 1000$ to those who get them. Given to the 20 who value them most,
they would be worth $20 \times (80 + 60)/2 = 1400$. The extra loss, 400, adds to the
triangle of 400 (if the lucky buyers cannot resell).
\end{solution}

\begin{solution}{exo:b1:price-theory-at-work:refute}
Yields that move far from $r + \delta - g$ for many years with no change in
depreciation or in any measure of expected growth; or yields that move
against the real rate when expected growth is steady, as measured by
surveys of house-price expectations. A single country's correlation does
not refute it, because $g$ is not observed.
\end{solution}

\begin{solution}{pb:b1:price-theory-at-work:store}
\textbf{1.} $200 \times 1.005 + 2 = 203$.

\textbf{2.} October 203, November 206.01, December 209.05, January 212.09,
February 215.15, March 218.23, April 221.32, May 224.42, June 227.55, July
230.68 (and August 233.84).

\textbf{3.} $230.68/200 - 1 = 15.34\%$.

\textbf{4.} Faster, and buying in September to sell later would earn more
than interest and storage, so more would be stored, raising the
September price and lowering later ones. Slower, and holders would sell
now and lend the money, which cannot happen while stocks are held.

\textbf{5.} $1.015^{10} - 1 = 16.05\%$.

\textbf{6.} The \omterm{def:b1:price-theory-at-work:carry}{cost of carry} for the month is $203 - 200 = 3$; the price rose 2. The
\omterm{def:b1:price-theory-at-work:carry}{convenience yield} is $1$ dalen a tonne.

\textbf{7.} It falls back to 200; stocks are run down to almost nothing just
before the harvest, since carrying old grain into a cheaper season loses
money.

\textbf{8.} Brenland 220 and Kestria 160.

\textbf{9.} $2(220 - p_K - 10) = 2(p_K - 160)$: $p_K = 185$, $p_B = 195$ and 50 thousand tonnes.

\textbf{10.} $195 = 185 + 10$: equality, since grain flows.

\textbf{11.} $2(220 - p_K - 30) = 2(p_K - 160)$: $p_K = 175$, $p_B = 205$, 30 thousand tonnes.

\textbf{12.} When $c \ge 220 - 160 = 60$.

\textbf{13.} Kestria's farmers gain (a higher price on more grain); Kestria's
buyers lose (a higher price).

\textbf{14.} $400 - 230.68 = 169.32$ thousand tonnes: the stock carried to July.

\textbf{15.} Demand at 210 is 190: a shortage of $190 - 169.32 = 20.68$.

\textbf{16.} The stock clears at a \omterm{def:b1:price-theory-at-work:full-price}{full price} of 230.68, so the queue costs
$230.68 - 210 = 20.68$ a tonne.

\textbf{17.} $20.68/10 = 2.07$ hours per tonne.

\textbf{18.} $20.68 \times 169.32 = 3501.5$ thousand dalens ($3502.1$ unrounded), about 3.5 million.

\textbf{19.} Carrying grain to July would no longer pay its cost, so less
would be stored: more sold early in the year, at lower prices, and less in
July, making the July shortage and queues worse.

\textbf{20.} $0.8 \times 10000 \times 12 = 96000$ dalens.

\textbf{21.} $96000/0.05 = 1.92$ million dalens.

\textbf{22.} 2.4 million at $4\%$ and 1.6 million at $6\%$.

\textbf{23.} $5\%$, the real rate (no depreciation, no growth).

\textbf{24.} Storage would be worth less, so the rent of storage space
would fall, and with it the store's value: the expected loss is
capitalised at once.

\textbf{25.} The store earns a lean-month premium of $15.34\%$: grain bought at
200 at harvest sells at 230.68 in July.
\end{solution}
